\chapter{Assembling a Driver: A Complete Worked Trace}
\label{ch:fulltrace}

\chapterorient{Every chapter so far has built one machine and set it on the shelf:
the mesh, the function spaces, the assembly fold, the weak forms, the Krylov
solvers, the multigrid preconditioners, the matrix-free applies, the reduction
verbs. This chapter takes the whole shelf down at once. We pick one concrete
simulation---the driven analysis of a two-port filter---and trace it from the
configuration file the engineer writes to the S-parameter curve the simulation
draws, naming, at every stage, the exact machine from the exact earlier chapter
that does the work. Nothing new is built here. Instead we watch \emph{all of it
run together}, see the whole simulator as a single composition of perhaps fifteen
named pieces, and---reading the same trace twice, once for the driven filter and
once for the simpler electrostatic capacitor---confirm the book's central claim
in motion: one simulation is one composition of the same handful of machines.
This is the synthesis of Parts~I through~V, and the bridge into Part~VI, where
that very composition is rendered as code.}

\section{One simulation, named all the way down}

We have spent five parts disassembling Palace. \Cref{part:modalities} took the
six analyses apart one at a time; before it, Part~IV built the numerical engines,
Part~III the calculus that names them, and Part~II the field theory they
discretize. Each chapter handed us a part and a name: \code{fe\_assemble},
\code{solve\_family}, \code{frequency\_sweep}, \code{ksp\_solve},
\code{gram\_reduce}, \code{sparameter\_reduce}, the V-cycle, the Chebyshev
smoother, the matrix-free apply. We have, by now, a complete vocabulary. What we
have \emph{not} yet done is run one real simulation slowly enough to point at
every part as it engages.

That is this chapter's only job. We will not derive a new equation or define a new
combinator. We will choose one problem, walk it through the lifecycle of
\cref{ch:lifecycle} stage by stage, and at each stage say: \emph{this is the
machine from that chapter, and here is what it does to the data right now}. The
value of the chapter is in its cross-references---it is the index of the whole
book's machinery caught in the act.

\begin{keyidea}
A complete Palace simulation is not a monolith. It is one \emph{composition} of a
small, fixed set of named machines---a dozen-and-a-half pieces, each built in an
earlier chapter, each with a precise type and a precise job. To ``understand a
driver'' is exactly to know which pieces it composes and in what order. By the end
of this chapter you will have seen the entire machine in a single expression, and
every box in that expression will be a chapter you have already read.
\end{keyidea}

\subsection{The problem we will trace}

We trace the analysis that exercises the most machinery: the \emph{driven
frequency sweep} of a two-port microwave \emph{filter} (\cref{ch:driven}). It is
the richest of the six because, end to end, it touches almost everything the book
built: it meshes a three-dimensional geometry, lays a Nédélec $\Hcurl$ space over
it, assembles three operators by the weak-form fold, forms a fresh complex
operator at every frequency, solves each by preconditioned GMRES wrapped around a
multigrid V-cycle whose smoother is Chebyshev/Hiptmair and whose every apply is
matrix-free, and finally projects the solved fields onto port modes borrowed from
a boundary-mode analysis. A single driven run is a guided tour of Parts~II--V.

Where the driven path grows too dense to read, we run the \emph{electrostatic
capacitance} extraction (\cref{ch:electrostatics}) beside it as the clean
contrast. It walks the identical lifecycle but lands in the simplest cells---the
fixed-operator map, the symmetric Gram---so it shows the same skeleton stripped to
its bones. Two analyses, one machine: the comparison is the chapter's spine.

\begin{physicsbox}
\textbf{The device.} An engineer hands us a two-port bandpass filter: a small
metal structure---coupled resonators on a substrate---with a coaxial or
waveguide \emph{port} at each end (\cref{fig:trace-pipeline}, far left). The
question is the network-analyzer question of \cref{ch:driven}: across a band of
frequencies $f \in [f_{\min}, f_{\max}]$, how much of a wave injected at port~1
is \emph{reflected} ($S_{11}$) and how much is \emph{transmitted} to port~2
($S_{21}$)? The answer is the scattering matrix $S(\omega)$ as a function of
frequency---the passband, the roll-off, the match. We will compute it by solving
the time-harmonic Maxwell field inside the filter at each frequency and reading
the scattered waves off the solved field.
\end{physicsbox}

\section{The lifecycle, walked once, naming every machine}
\label{sec:trace-walk}

Recall the lifecycle of \cref{ch:lifecycle}: parse, configure, dispatch, build the
mesh, solve, (optionally refine), output. We now walk those stages \emph{for this
one run}, and at each stage we name the component---and the chapter---doing the
work. \Cref{fig:trace-pipeline} is the whole walk in one annotated picture; keep
it in view, because the subsections below are simply its boxes, read left to right.

\begin{figure}[p]
\centering
\begin{tikzpicture}[node distance=6mm and 6mm, font=\footnotesize, >=Stealth]
  % the spine, left to right
  \node[data, align=center, text width=17mm] (cfg)
    {config file\\[1pt]\scriptsize geometry,\\$\varepsilon,\mu$, ports,\\freq.\ band};
  \node[stage, right=of cfg, align=center, text width=20mm] (io)
    {\textbf{parse}\\[1pt]\scriptsize \code{IoData}\\record};
  \node[stage, right=of io, align=center, text width=18mm] (mesh)
    {\textbf{mesh}\\[1pt]\scriptsize \code{build\_mesh}};
  \node[stage, right=of mesh, align=center, text width=18mm, fill=simBgViolet, draw=simViolet] (disp)
    {\textbf{dispatch}\\[1pt]\scriptsize driven driver};
  \node[stage, below=14mm of io, align=center, text width=21mm] (asm)
    {\textbf{assemble}\\[1pt]\scriptsize \code{fe\_assemble}$^{\times3}$\\$\{\Kop,\Cop,\Mop\}$};
  \node[stage, right=of asm, align=center, text width=24mm, fill=simBgRust, draw=simRust] (solve)
    {\textbf{solve}\\[1pt]\scriptsize \code{frequency\_sweep}\\(per $\omega$: GMRES)};
  \node[stage, right=of solve, align=center, text width=21mm, fill=simBgGold, draw=simGold] (red)
    {\textbf{reduce}\\[1pt]\scriptsize \code{sparameter\_}\\\code{reduce}};
  \node[product, right=of red, align=center, text width=15mm] (prod)
    {$S(\omega)$\\[1pt]\scriptsize S-param\\plot};
  % top-row forward flow
  \draw[flow] (cfg) -- (io);
  \draw[flow] (io) -- (mesh);
  \draw[flow] (mesh) -- (disp);
  % down into the solve band
  \draw[flow] (disp.south) |- (asm.north west);
  % the solve band, left to right
  \draw[flow] (asm) -- node[above,font=\scriptsize]{\code{fam}} (solve);
  \draw[flow] (solve) -- node[above,font=\scriptsize]{$[\vE_k]$} (red);
  \draw[flow] (red) -- (prod);
  % chapter tags under each box
  \node[cornertag, below=1mm of io]   {\cref{ch:records}};
  \node[cornertag, below=1mm of mesh] {\cref{ch:functionspaces}};
  \node[cornertag, below=1mm of disp] {\cref{ch:lifecycle}};
  \node[cornertag, below=1mm of asm]  {\cref{ch:assembly}, \cref{ch:weak}};
  \node[cornertag, below=1mm of solve]{\cref{ch:gmres}, \cref{ch:multigrid}};
  \node[cornertag, below=1mm of red]  {\cref{ch:reductions}, \cref{ch:waveports}};
  % the port modes feeding the reduction from the side
  \node[extern, above=8mm of red, align=center, text width=20mm] (ports)
    {port modes $\vs_i$\\[1pt]\scriptsize boundary-mode\\(\cref{ch:waveports})};
  \draw[refedge] (ports.south) -- (red.north);
\end{tikzpicture}
\caption[The complete driven pipeline, annotated by chapter]{The capstone
diagram: one complete driven simulation, from the configuration file to the
S-parameter plot, with every stage tagged by the chapter that built it. The
top row is the shared lifecycle set-up (parse the \code{IoData} record, build the
mesh, dispatch on the problem type, \cref{ch:lifecycle}); the bottom band is the
driver's \texttt{assemble--solve--reduce} body (\cref{ch:situating}). The
\emph{assemble} box folds three weak-form term lists into the fixed basis
$\{\Kop,\Cop,\Mop\}$ (\cref{ch:assembly}); the \emph{solve} box is the
operator-varying frequency sweep, each frequency a preconditioned GMRES solve
(\cref{ch:gmres}) over a multigrid V-cycle (\cref{ch:multigrid}); the
\emph{reduce} box projects the solved fields onto the port modes that the
boundary-mode analysis supplies from the side (\cref{ch:waveports}). Every box is
a machine from an earlier part; this chapter watches them all run together.}
\label{fig:trace-pipeline}
\end{figure}

\subsection{Parse the configuration: the \texorpdfstring{\code{IoData}}{IoData} record}
\label{sec:trace-parse}

The run begins with one file the engineer wrote: a configuration naming the mesh,
the polynomial order, the materials ($\varepsilon$ and $\mu$ per region), the two
ports and their excitations, the frequency band to sweep, and---the field that
decides everything downstream---the problem type, \code{Driven}. Parsing it
produces the single in-memory value the rest of the run consults: the
\code{config}, which is the \code{IoData} record of \cref{ch:records}.

The decisive property, established in \cref{ch:lifecycle} and \cref{ch:strata}, is
that this record is \emph{read-only}. It is built once and never modified; every
later stage reads it, none writes it. It is \emph{construction-stratum} data
(\cref{ch:strata}): fixed the instant parsing finishes, so any computation that
depends only on it is referentially stable for the whole run. That single fact is
what licenses every hoist, reuse, and reorder we are about to watch.

\begin{notationbox}
We write \code{config} for the parsed \code{IoData} value. It is a record
(\cref{ch:records})---a bundle of named fields read with a dot:
\lstinline[style=l4style]|config.problem_type|, the materials
\lstinline[style=l4style]|config.materials|, the swept band
\lstinline[style=l4style]|config.frequencies|, the ports
\lstinline[style=l4style]|config.ports|. The configuration's records get their
definitional home in \cref{ch:records}; here we simply read fields off it. The
two analyses we trace read the \emph{same} record type and differ only in which
fields they consult---the electrostatic path reads \code{config.terminals} where
the driven path reads \code{config.frequencies} and \code{config.ports}.
\end{notationbox}

\subsection{Configure the device, then dispatch}

Two short stages follow, both from \cref{ch:lifecycle}. \emph{Configure the
device} brings up the CPU or GPU backend the element kernels will run on---the
choice read from \code{config}, the computation downstream expressed against an
abstract device so the same description runs on either. Then \emph{dispatch}: the
run reads \code{config.problem\_type}, finds \code{Driven}, and on the strength of
that one value selects the \emph{driven driver}---the object that knows how to run
this one analysis end to end. This is the \emph{specialization seam} of
\cref{ch:lifecycle}: everything above it is shared by all six analyses; below it,
the driven driver fills in the one stage that differs. Our electrostatic contrast
takes the same seam to the \emph{electrostatic} driver instead---one field, one
\code{switch}, the entire difference of \emph{which} analysis runs.

\subsection{Build the mesh, and lay the \texorpdfstring{$\Hcurl$}{H(curl)} space}
\label{sec:trace-mesh}

With a driver chosen, the run builds the computational mesh: it loads the geometry
named in \code{config}, applies any preprocessing, partitions it, and performs any
a-priori refinement requested. The framework verb is \code{build\_mesh}, and the
value it returns---the \code{Mesh}---is the first value the lifecycle
\emph{produces} rather than reads (the config came from the user; the mesh is
computed). Mesh construction is driver-agnostic in shape: both our analyses build
the mesh identically.

Over that mesh the driver lays its \emph{finite-element space} (\cref{ch:functionspaces}).
Here the two analyses first diverge in physics. The driven analysis solves for a
vector field $\vE$ that must support tangential continuity and a curl, so it lives
in the Nédélec $\Hcurl$ space---the edge-element slot of the de Rham complex
(\cref{ch:derham}). The electrostatic analysis solves for a scalar potential $V$,
so it lives in the nodal $\Hone$ space---the vertex slot. The choice of space is
\emph{one cell} of the grid of \cref{ch:situating}, and it is the first cell our
two traces fill differently: $\Hcurl$ for driven, $\Hone$ for electrostatic.

\begin{remark}[The adaptive loop, deferred]
The lifecycle wraps the solve in the adaptive estimate--mark--refine loop of
\cref{ch:amr}---the \emph{outer fold} whose threaded value is the mesh. For our
trace we keep refinement off, so the fold runs its body exactly once and the
lifecycle is the straight line of \cref{ch:lifecycle}. Everything below is
\emph{one pass} of that fold. With refinement on, the entire
assemble--solve--reduce body we are about to walk simply runs again on a finer
mesh, each pass beginning from the mesh the last produced---the
\code{fold\_solve} of \cref{ch:fold} wrapped around all of it.
\end{remark}

\subsection{Assemble the basis, once: the weak-form fold}
\label{sec:trace-assemble}

Now the driver's body begins, and the first stage is \emph{assemble}. The driven
analysis needs three operators on $\Hcurl$, the same three characters from every
chapter: the curl--curl stiffness $\Kop$ (built from the reluctivity $\nu=1/\mu$),
the damping $\Cop$ (from surface impedance / conductivity), and the mass $\Mop$
(from permittivity $\varepsilon$). Each is the result of the assembly fold
\code{fe\_assemble} of \cref{ch:assembly}, summing one contribution per weak-form
term (\cref{ch:weak}) over the mesh, with the essential boundary degrees of
freedom eliminated (\cref{ch:bc}):
\begin{lstlisting}[style=l4style]
-- the fixed basis: three weak-form folds, run ONCE, before any sweep
space = nd_space cfg                                     -- H(curl) Nedelec space
fam   = { K  = fe_assemble space [ curl_curl (reluctivity cfg) ]   -- stiffness
        , C  = fe_assemble space [ damping cfg ]                   -- damping
        , M  = fe_assemble space [ mass (permittivity cfg) ] }     -- mass
\end{lstlisting}
Three structural facts, each a load-bearing reference. First, every
\code{fe\_assemble} is the \emph{fold} of \cref{ch:assembly}---a loop over the term
list that we do not write; the framework folds it. Second, each operator's matrix
action, when the solver later applies it, is \emph{matrix-free}
(\cref{ch:matrixfree}): the operator is never stored as an explicit sparse matrix
but applied as a sequence of element-local tensor contractions, the libCEED-style
quadrature kernel. Third---and this is the word from \cref{ch:electrostatics} that
governs the whole solve---the basis is assembled \emph{once}, before the sweep, and
held. The electrostatic contrast is the same fold run once over a \emph{single}
term, the $\varepsilon$-weighted gradient, producing one operator $\Kop$.

\begin{keyidea}
The \emph{assemble} stage is the weak-form fold \code{fe\_assemble}
(\cref{ch:assembly}), run once per operator before any solving. The driven
analysis folds three term lists into a fixed basis $\{\Kop,\Cop,\Mop\}$; the
electrostatic analysis folds one into $\Kop$. The operators' applies are
matrix-free (\cref{ch:matrixfree}); the essential boundary data is eliminated
(\cref{ch:bc}); the result is read-only and reused across the entire solve. ``Once''
is the structural signature shared by both analyses---it is the upstream half of
the hoist.
\end{keyidea}

\subsection{Solve: the operator-varying sweep, and its nested engine}
\label{sec:trace-solve}

The solve stage is where the driven analysis spends almost all of its time, and
where the most machinery engages at once. Recall from \cref{ch:driven} that the
system operator is \emph{not} fixed: at each frequency $\omega$ in the swept band
we form a fresh complex operator from the fixed basis,
\begin{equation}
  \mat{A}(\omega) \;=\; \Kop \;+\; i\omega\,\Cop \;-\; \omega^2\Mop,
  \label{eq:trace-Aomega}
\end{equation}
solve $\mat{A}(\omega)\,\vE = \vb(\omega)$ for the phasor field, and move on. The
solves are \emph{independent}---the field at one frequency needs nothing from
another---so the sweep is a \emph{map}, the operator-varying map (corner~2) of
\cref{ch:situating}. The framework verb is \code{frequency\_sweep}, and forming
\cref{eq:trace-Aomega} from the basis is the \code{linear\_combination} combinator
of \cref{ch:combinators}, specialized to operator operands and complex scalars.

What makes this stage worth a figure of its own is that each single
\code{ksp\_solve} inside the sweep is itself a \emph{stack} of the machines from
Part~IV, nested four deep (\cref{fig:trace-solverstack}):

\begin{enumerate}[nosep]
  \item \textbf{GMRES} (\cref{ch:gmres}) drives the outer iteration. The operator
  $\mat{A}(\omega)$ is complex, non-Hermitian, and indefinite---the $i\omega\Cop$
  term makes it complex, the $-\omega^2\Mop$ term makes it indefinite---so none of
  CG's structure survives and GMRES is the right Krylov engine (\cref{ch:cg}
  versus \cref{ch:gmres}).
  \item \textbf{A multigrid V-cycle} (\cref{ch:multigrid}) is the
  \emph{preconditioner} GMRES calls each iteration: restrict the residual down the
  mesh hierarchy, solve cheaply on the coarsest level, prolong the correction back
  up.
  \item \textbf{A Chebyshev / Hiptmair smoother} (\cref{ch:smoothers}) runs at
  each level of the V-cycle, damping the high-frequency error the coarse grid
  cannot see---the Hiptmair variant handling the $\Hcurl$ null-space that an
  ordinary smoother would miss.
  \item \textbf{A matrix-free apply} (\cref{ch:matrixfree}) is the innermost
  operation of all of them. Every GMRES Arnoldi step, every smoother sweep, every
  residual evaluation ultimately \emph{applies $\mat{A}(\omega)$ to a vector}, and
  that apply is the element-local tensor contraction of \cref{ch:matrixfree}---no
  assembled sparse matrix ever materializes.
\end{enumerate}

Threaded through all four levels is the solver \emph{state}---the Krylov basis, the
running residual, the iteration counters---carried by the solve monad of
\cref{ch:monad} as the \code{SimState} record (\cref{ch:records}), so that the
mutation the C++ performs in place is, in our reading, a pure state thread.

\begin{figure}[p]
\centering
\begin{tikzpicture}[font=\footnotesize, >=Stealth]
  % nested boxes: GMRES > V-cycle > smoother > matrix-free apply
  \node[stage, minimum width=72mm, minimum height=52mm, fill=simBgRust,
        draw=simRust, align=center] (g) at (0,0) {};
  \node[anchor=north, font=\small\bfseries, text=simRust] at (g.north)
    {\ GMRES (\cref{ch:gmres})};
  \node[font=\scriptsize\itshape, text=simRust, anchor=north]
    at ($(g.north)+(0,-0.5)$) {outer Krylov: complex, non-Hermitian $\mat{A}(\omega)$};

  \node[stage, minimum width=60mm, minimum height=36mm, fill=simBgBlue,
        draw=simBlue, align=center] (v) at (0,-0.55) {};
  \node[anchor=north, font=\footnotesize\bfseries, text=simBlue] at (v.north)
    {\ multigrid V-cycle (\cref{ch:multigrid})};
  \node[font=\scriptsize\itshape, text=simBlue, anchor=north]
    at ($(v.north)+(0,-0.42)$) {preconditioner: restrict $\to$ coarse-solve $\to$ prolong};

  \node[stage, minimum width=46mm, minimum height=20mm, fill=simBgGreen,
        draw=simGreen, align=center] (s) at (0,-1.1) {};
  \node[anchor=north, font=\footnotesize\bfseries, text=simGreen] at (s.north)
    {\ Chebyshev / Hiptmair smoother};
  \node[font=\scriptsize\itshape, text=simGreen, anchor=north]
    at ($(s.north)+(0,-0.4)$) {(\cref{ch:smoothers}): damp high-freq.\ error};

  \node[extern, minimum width=32mm, minimum height=7mm] (mf) at (0,-1.75)
    {matrix-free apply (\cref{ch:matrixfree})};

  % a side label: the threaded state
  \node[opbox, anchor=west, align=center, text width=24mm, fill=white]
    (state) at (4.6,1.4) {\code{SimState}\\\scriptsize threaded by the\\solve monad\\(\cref{ch:monad}, \cref{ch:records})};
  \draw[refedge] (state.south) .. controls +(0,-1.4) and +(0,1.2) .. (g.east);
  % apply call-down annotation
  \node[font=\scriptsize\itshape, text=simViolet, anchor=west] at (2.7,-1.75)
    {$\leftarrow$ every level};
  \node[font=\scriptsize\itshape, text=simViolet, anchor=west] at (2.7,-2.15)
    {\ \ bottoms out here};
\end{tikzpicture}
\caption[The solver stack for one frequency]{The solver stack engaged by a
\emph{single} frequency of the sweep---four machines from Part~IV, nested. The
outermost is preconditioned \textbf{GMRES} (\cref{ch:gmres}), the only Krylov
method that handles the complex, non-Hermitian, indefinite $\mat{A}(\omega)$. Its
preconditioner is a \textbf{multigrid V-cycle} (\cref{ch:multigrid}), which at each
mesh level calls a \textbf{Chebyshev/Hiptmair smoother} (\cref{ch:smoothers}) to
damp the error the coarse grid cannot resolve. Beneath all of them, every
operator application---every Arnoldi step, every smoother sweep, every residual---is
a \textbf{matrix-free apply} (\cref{ch:matrixfree}): an element-local tensor
contraction, never an assembled sparse matrix. The whole stack's mutable state (the
Krylov basis, residual, counters) is the \code{SimState} record threaded by the
solve monad (\cref{ch:monad}). And this entire stack runs \emph{once per
frequency}, fresh operator each time.}
\label{fig:trace-solverstack}
\end{figure}

The electrostatic contrast collapses this stack dramatically, and the collapse is
instructive. Its operator $\Kop$ is symmetric positive-definite and \emph{fixed},
so the outer Krylov method is \emph{CG}, not GMRES (\cref{ch:cg}); the
preconditioner is built \emph{once} and reused across the whole family (the hoist
of \cref{ch:electrostatics}); and the sweep is the fixed-operator map
\code{solve\_family} (corner~1) rather than \code{frequency\_sweep} (corner~2). The
inner machines---a V-cycle, a smoother, a matrix-free apply---are the \emph{same
machines}; only the outer Krylov method and the position of the operator capture
change. \emph{One stack, two configurations.}

\begin{pitfall}
The single line whose position separates the two analyses is the operator capture
inside \code{frequency\_sweep}. In electrostatics it is \emph{hoisted} above the
loop (the operator never changes); in the driven sweep it lives \emph{inside} the
loop (the operator changes every frequency). Moving it the wrong way in the driven
sweep does not slow the code---it solves every frequency against
$\mat{A}(\omega_0)$, the first frequency's operator, and returns a smooth,
plausible, \emph{wrong} curve (\cref{ch:driven}). The loop position of one capture
call is part of the algorithm, not a performance knob.
\end{pitfall}

\subsection{Reduce: the port-projection, with borrowed modes}
\label{sec:trace-reduce}

The sweep hands the reduction a family of complex fields $[\vE_k]$, one per swept
frequency. The reduction collapses each into the scattering parameters. This is
the rank-2 \emph{port-projection} of \cref{ch:reductions}---reduction shape~2 of
\cref{ch:situating}, and emphatically \emph{not} a Gram. The verb is
\code{sparameter\_reduce}, applied once per frequency:
\begin{equation}
  S_{ij}(\omega) \;=\; \sigma_{ij}\,\bigl(\inner{\vs_i}{\vE_j(\omega)} - \delta_{ij}\bigr),
  \label{eq:trace-sparam}
\end{equation}
the projection of the drive-$j$ field onto receiver mode $i$, minus the Kronecker
self-term on the diagonal (which converts ``total field at the driven port'' into
``reflected field''), scaled by the port-kind factor $\sigma_{ij}$.

The one component this stage \emph{consumes from elsewhere} is the set of port
modes $\vs_i$. They do not come from the driven solve; they come from the
\emph{boundary-mode analysis} of \cref{ch:waveports}, which solves a small
two-dimensional eigenproblem on each port's cross-section to find the field
pattern that propagates there. In the pipeline diagram
(\cref{fig:trace-pipeline}) this is the arrow entering the reduce box from the
side: the driven driver \emph{composes} the boundary-mode driver's product as one
of its inputs. The electrostatic contrast instead reduces with the symmetric Gram
\code{gram\_reduce} (\cref{ch:reductions}), $C_{ij}=V_j^{\!\top}\Kop V_i$, pairing
the solved family \emph{with itself} through the energy operator---no borrowed
basis, symmetric by construction.

\subsection{Output}

When the reduction finishes, the driver writes its product---the table of
$S_{ij}(\omega)$ across the band---to disk, and the top-level lifecycle writes the
run metadata around it (\cref{ch:lifecycle}). Plotted, $S(\omega)$ is the two
familiar microwave curves: the return loss $\abs{S_{11}}$ in dB and the insertion
loss $\abs{S_{21}}$ in dB (\cref{fig:trace-splot}). That plot is the answer the
engineer ran the simulation to obtain.

\section{The whole machine in one expression}
\label{sec:trace-composition}

We have walked the stages. Now collapse the walk into a single composition and
\emph{expand it}, so the reader sees the entire machine---every nested engine---in
one tree.

\subsection{The three-stage composition}

At the top level the driven driver is the three-stage composition of
\cref{ch:driven}:
\begin{lstlisting}[style=l4style]
driven = sparameter_reduce . frequency_sweep . fe_assemble
\end{lstlisting}
read right to left: assemble the basis once, sweep the operator-varying map over
the band, reduce by port-projection. Spelled out with its data, exactly the
composition of \cref{ch:driven}:
\begin{lstlisting}[style=l4style]
-- inputs = config; output = the frequency response S(omega)
driven :: DrivenConfig -> FrequencyResponse
driven cfg =
  let space  = nd_space cfg                                       -- H(curl) space
      fam    = { K = fe_assemble space [ curl_curl (reluctivity cfg) ]  -- (1) assemble
               , C = fe_assemble space [ damping cfg ]                  --     the basis
               , M = fe_assemble space [ mass (permittivity cfg) ] }    --     ONCE
      omegas = sample_frequencies cfg                             -- the swept band
      es     = frequency_sweep fam omegas       -- (2) operator-VARYING map: A(w) rebuilt inside
  in  sparameter_reduce (ports cfg) es          -- (3) port-projection -> S(omega)
\end{lstlisting}

\subsection{Expanded all the way down}

The three boxes hide the whole of Part~IV. Expand each piece and the entire call
tree appears---the assemble fold, the per-$\omega$ solve with its full
preconditioner stack, the reduction over borrowed modes. This is the hero figure
of the chapter: the whole simulator in one tree (\cref{fig:trace-tree}).

\begin{lstlisting}[style=l4style]
driven cfg
|
+- (1) fe_assemble x3              -- the basis, folded ONCE   [ch:assembly, ch:weak]
|     +- fold over weak-form terms                            -- [ch:assembly]
|     +- element-local quadrature  -- matrix-free contraction  [ch:matrixfree]
|     +- eliminate essential BC dofs                          -- [ch:bc]
|
+- (2) frequency_sweep fam omegas  -- map over the band (corner 2) [ch:driven]
|     +- FOR EACH omega:           -- independent: a map, not a fold
|         +- A(w) = linear_combination [(1,K),(iw,C),(-w^2,M)]  -- [ch:combinators]
|         +- b    = excitation fam omega                       -- per-omega RHS
|         +- ksp_solve A(w) b      -- one GMRES solve           [ch:gmres]
|             +- GMRES outer iteration                          -- [ch:gmres]
|             |   +- precondition: multigrid V-cycle            -- [ch:multigrid]
|             |   |   +- restrict / coarse-solve / prolong
|             |   |   +- smoother (Chebyshev / Hiptmair)        -- [ch:smoothers]
|             |   |       +- matrix-free apply of A(w)          -- [ch:matrixfree]
|             |   +- Arnoldi step: matrix-free apply of A(w)    -- [ch:matrixfree]
|             +- state threaded by the solve monad (SimState)  -- [ch:monad, ch:records]
|
+- (3) sparameter_reduce ports es  -- port-projection (shape 2) [ch:reductions]
      +- S_ij = scale * (project s_i E_j - selfterm i j)       -- [ch:reductions]
      +- port modes s_i FROM boundary-mode analysis            -- [ch:waveports]
\end{lstlisting}

\begin{figure}[p]
\centering
\begin{tikzpicture}[font=\scriptsize, >=Stealth,
    grow=down, level distance=8mm,
    edge from parent/.style={draw=simGray, semithick, -{Stealth[length=1.6mm]}},
    every node/.style={align=center}]
  \node[stage, fill=simBgViolet, draw=simViolet, text width=26mm]
    {\code{driven cfg}\\\scriptsize the whole machine}
    child {node[stage, text width=24mm] {\textbf{(1)} \code{fe\_assemble}$^{\times3}$\\\scriptsize basis, once}
      child {node[opbox, text width=22mm] {fold over\\weak-form terms\\\scriptsize\cref{ch:assembly}}}
      child {node[extern, text width=20mm] {matrix-free\\quadrature\\\scriptsize\cref{ch:matrixfree}}}
    }
    child {node[stage, fill=simBgRust, draw=simRust, text width=30mm] {\textbf{(2)} \code{frequency\_sweep}\\\scriptsize map over $\omega$ (corner 2)}
      child {node[opbox, text width=26mm] {$\mat{A}(\omega)=$\code{lin\_comb}\\\scriptsize\cref{ch:combinators}}}
      child {node[stage, fill=simBgRust, draw=simRust, text width=22mm] {\code{ksp\_solve}\\\scriptsize GMRES \cref{ch:gmres}}
        child {node[stage, fill=simBgBlue, draw=simBlue, text width=22mm] {V-cycle\\precond.\\\scriptsize\cref{ch:multigrid}}
          child {node[stage, fill=simBgGreen, draw=simGreen, text width=20mm] {smoother\\\scriptsize\cref{ch:smoothers}}
            child {node[extern, text width=20mm] {matrix-free\\apply\\\scriptsize\cref{ch:matrixfree}}}
          }
        }
      }
    }
    child {node[stage, fill=simBgGold, draw=simGold, text width=26mm] {\textbf{(3)} \code{sparameter\_reduce}\\\scriptsize projection (shape 2)}
      child {node[opbox, text width=22mm] {$S_{ij}$ projection\\\scriptsize\cref{ch:reductions}}}
      child {node[extern, text width=20mm] {port modes $\vs_i$\\\scriptsize\cref{ch:waveports}}}
    };
\end{tikzpicture}
\caption[The driven driver, expanded to its leaves]{The hero figure: the entire
driven driver expanded from the top-level composition down to its matrix-free
leaves. The three top-level stages---\code{fe\_assemble} (assemble the basis once),
\code{frequency\_sweep} (the operator-varying solve map), and
\code{sparameter\_reduce} (the port-projection)---each unfold into the machines
that implement them. The solve branch is the deepest: a GMRES outer iteration
(\cref{ch:gmres}) preconditioned by a multigrid V-cycle (\cref{ch:multigrid})
whose smoother (\cref{ch:smoothers}) bottoms out, like every operator application
in the tree, in a matrix-free apply (\cref{ch:matrixfree}). Read the whole tree
and you have read the whole book: every leaf is a chapter. This is exactly the
``synthesized library'' view that Part~VI (\cref{ch:synthlib}, \cref{ch:composing})
renders as code.}
\label{fig:trace-tree}
\end{figure}

\begin{keyidea}
The whole simulator is \emph{one expression}. \texttt{driven $=$
sparameter\_reduce $\circ$ frequency\_sweep $\circ$ fe\_assemble}, and each of
those three pieces expands into the machines of Parts~III--IV until every leaf is a
named component from an earlier chapter. There is no hidden glue, no special-case
plumbing: a driver \emph{is} its composition tree, and the tree is finite and
shallow. You have now seen the entire machine in one figure---and you have read
every box in it.
\end{keyidea}

\section{Two analyses, one machine}
\label{sec:trace-twoanalyses}

The driven trace is one path through the machinery. The chapter's claim is that the
\emph{same} machinery, recomposed, is every other analysis too. We make that
concrete by setting the electrostatic capacitance trace beside the driven filter
trace and reading off what is \emph{shared} and what is \emph{swapped}
(\cref{fig:trace-contrast}).

\begin{figure}[t]
\centering
\begin{tikzpicture}[font=\footnotesize, >=Stealth, node distance=5mm and 7mm]
  % ===== DRIVEN row =====
  \node[font=\small\bfseries, text=simRust, anchor=west] at (-0.3,2.0) {Driven filter};
  \node[stage, align=center, text width=15mm] (dasm) at (1.4,1.4) {assemble\\\scriptsize$\{\Kop,\Cop,\Mop\}$};
  \node[stage, fill=simBgRust, draw=simRust, align=center, text width=18mm] (dsolve) at (4.0,1.4)
    {\code{frequency\_}\\\code{sweep}\\\scriptsize GMRES, corner~2};
  \node[stage, fill=simBgGold, draw=simGold, align=center, text width=17mm] (dred) at (6.7,1.4)
    {\code{sparameter\_}\\\code{reduce}};
  \node[product, align=center, text width=11mm] (dprod) at (8.9,1.4) {$S(\omega)$};
  \draw[flow] (dasm) -- (dsolve); \draw[flow] (dsolve) -- (dred); \draw[flow] (dred) -- (dprod);
  % ===== ELECTROSTATIC row =====
  \node[font=\small\bfseries, text=simBlue, anchor=west] at (-0.3,-0.3) {Electrostatic capacitor};
  \node[stage, align=center, text width=15mm] (easm) at (1.4,-0.9) {assemble\\\scriptsize$\Kop$ only};
  \node[stage, fill=simBgBlue, draw=simBlue, align=center, text width=18mm] (esolve) at (4.0,-0.9)
    {\code{solve\_family}\\\scriptsize CG, corner~1};
  \node[stage, fill=simBgGold, draw=simGold, align=center, text width=17mm] (ered) at (6.7,-0.9)
    {\code{gram\_reduce}\\\scriptsize $w{=}1$};
  \node[product, align=center, text width=11mm] (eprod) at (8.9,-0.9) {$\mat{C}$};
  \draw[flow] (easm) -- (esolve); \draw[flow] (esolve) -- (ered); \draw[flow] (ered) -- (eprod);
  % ===== shared / swapped annotations =====
  \node[font=\scriptsize\itshape, text=simGreen, align=center] at (1.4,0.25)
    {\textbf{shared:}\\\code{fe\_assemble} fold\\(\cref{ch:assembly})};
  \node[font=\scriptsize\itshape, text=simRust, align=center] at (4.0,0.25)
    {\textbf{swapped:}\\corner~1 $\leftrightarrow$ 2\\CG $\leftrightarrow$ GMRES};
  \node[font=\scriptsize\itshape, text=simRust, align=center] at (6.7,0.25)
    {\textbf{swapped:}\\Gram $\leftrightarrow$ projection};
  % shared inner-stack note
  \node[font=\scriptsize\itshape, text=simGreen, align=center] at (4.0,-2.1)
    {\textbf{shared inner stack:} V-cycle (\cref{ch:multigrid}) $\cdot$ smoother (\cref{ch:smoothers}) $\cdot$ matrix-free apply (\cref{ch:matrixfree})};
\end{tikzpicture}
\caption[Two analyses, one machine]{The driven filter and the electrostatic
capacitor as two recompositions of one machine. \textbf{Shared} (green): both
build their operators with the \emph{same} assembly fold \code{fe\_assemble}
(\cref{ch:assembly}), and both solves bottom out in the \emph{same} inner
stack---a multigrid V-cycle (\cref{ch:multigrid}) over a smoother
(\cref{ch:smoothers}) over a matrix-free apply (\cref{ch:matrixfree}).
\textbf{Swapped} (rust): the solve \emph{corner} (fixed-operator map, corner~1,
versus operator-varying map, corner~2) and with it the outer Krylov method (CG
versus GMRES); and the \emph{reduction} (symmetric Gram versus port-projection).
Two boxes move; the skeleton, and most of the inner machinery, is identical. This
is the claim of \cref{ch:situating} caught in a single picture.}
\label{fig:trace-contrast}
\end{figure}

What is \emph{shared} is most of the machine. Both analyses parse the same
\code{IoData} record, build the mesh the same way, dispatch through the same seam,
assemble with the same fold \code{fe\_assemble}, and---this is the part worth
dwelling on---bottom out in the \emph{same} inner solver stack: a multigrid
V-cycle over a smoother over a matrix-free apply. The expensive numerical core is
identical; the analyses do not each carry their own private solver.

What is \emph{swapped} is small and well-localized: the function space ($\Hcurl$
versus $\Hone$), the solve corner (operator-varying map versus fixed-operator map,
and with it GMRES versus CG), and the reduction (port-projection versus symmetric
Gram). This is precisely the ``one stage swapped'' punchline of
\cref{ch:situating} and \cref{ch:electrostatics}, now visible as two real traces
side by side rather than a table of dashes.

\begin{keyidea}
The driven filter and the electrostatic capacitor are the \emph{same machine} in
two configurations. They share the lifecycle, the \code{IoData} record, the mesh
build, the assembly fold, and---crucially---the entire inner solver stack
(V-cycle, smoother, matrix-free apply). They differ in exactly three swapped
cells: the function space, the solve corner (with its Krylov method), and the
reduction. A different analysis again would swap a different few cells and reuse
everything else. The simulator is not six programs; it is one composition with a
handful of swappable parts.
\end{keyidea}

\section{Counting the corners and shapes this one run touched}
\label{sec:trace-count}

Step back to the taxonomy of \cref{ch:situating} and ask: of the four solve
corners and three reduction shapes, how many did this \emph{single} driven run
use? The answer makes the economy of the design tangible (\cref{fig:trace-grid}).

The driven run touched \emph{one} solve corner---the operator-varying map
(corner~2)---and \emph{one} reduction shape---the rank-2 port-projection (shape~2).
It used, inside that one corner, a stack of Krylov and multigrid layers, and it
\emph{borrowed} one product from another analysis (the port modes from the
black-box boundary-mode corner, corner~4). That is the entire taxonomic footprint
of a complete microwave simulation: one corner, one reduction, a shared inner
stack, one borrowed input.

\begin{figure}[t]
\centering
\begin{tikzpicture}[font=\footnotesize]
  % the 2x2 corners grid + the table, with THIS run's cells highlighted
  \node[cornertag, anchor=south] at (0.9,2.5) {operator \emph{fixed}};
  \node[cornertag, anchor=south] at (3.5,2.5) {operator \emph{varying}};
  \node[cornertag, rotate=90, anchor=south] at (-1.4,1.8) {\textbf{map}};
  \node[cornertag, rotate=90, anchor=south] at (-1.4,0.4) {\textbf{fold}};
  % top-left: fixed map (electrostatic uses; driven does NOT)
  \node[stage, minimum width=20mm, minimum height=12mm, fill=simBgGray, draw=simGray,
        align=center] (c1) at (0.9,1.7) {\scriptsize fixed-op map\\\scriptsize(electrostatic)};
  % top-right: operator-varying map -- THIS RUN
  \node[stage, minimum width=20mm, minimum height=12mm, fill=simBgRust, draw=simRust,
        very thick, align=center] (c2) at (3.5,1.7) {\scriptsize \textbf{op-varying map}\\\scriptsize\textbf{(driven: HERE)}};
  % bottom-left: fold
  \node[stage, minimum width=20mm, minimum height=12mm, fill=simBgGray, draw=simGray,
        align=center] (c3) at (0.9,0.3) {\scriptsize state fold\\\scriptsize(transient/AMR)};
  % bottom-right: black box
  \node[extern, minimum width=20mm, minimum height=12mm, align=center]
        (c4) at (3.5,0.3) {\scriptsize black-box\\\scriptsize(borrowed: ports)};
  % the borrow arrow
  \draw[refedge, simViolet] (c4.north) -- node[right, font=\scriptsize, text=simViolet]{modes} (c2.south);
  % the reductions column, to the right
  \node[font=\scriptsize\bfseries, anchor=south] at (7.2,2.5) {reduction shape};
  \node[stage, minimum width=24mm, minimum height=7mm, fill=simBgGray, draw=simGray]
        (r1) at (7.2,1.75) {\scriptsize Gram (electrostatic)};
  \node[stage, minimum width=24mm, minimum height=7mm, fill=simBgGold, draw=simGold,
        very thick] (r2) at (7.2,0.95) {\scriptsize \textbf{projection (driven)}};
  \node[stage, minimum width=24mm, minimum height=7mm, fill=simBgGray, draw=simGray]
        (r3) at (7.2,0.15) {\scriptsize table ($f,Q$/energy)};
\end{tikzpicture}
\caption[The taxonomic footprint of one run]{What one complete driven run touches
on the grid of \cref{ch:situating}. Of the four solve corners (left $2\times2$),
the driven filter occupies exactly \emph{one}: the operator-varying map (corner~2,
rust). Of the three reduction shapes (right column), it uses exactly \emph{one}:
the rank-2 port-projection (gold). It \emph{borrows} the port modes from the
black-box corner's boundary-mode analysis (violet arrow). The electrostatic
contrast (gray) occupies the fixed-operator map and the Gram instead. A complete
microwave simulation is one corner, one reduction, a shared inner stack, and one
borrowed input---and a different analysis swaps only which cells light up. The grid
is small because the simulator is: a few corners $\times$ a few reductions $\times$
the shared field theory and calculus.}
\label{fig:trace-grid}
\end{figure}

This is the payoff of the whole taxonomy. A simulator that looked, from outside, like
six elaborate programs is, from inside, a \emph{few corners} $\times$ a \emph{few
reductions}, all resting on one shared field theory (Part~II) and one shared
calculus (Part~III). The driven run lit up one cell of each axis. The electrostatic
contrast lit up a different cell of each. Every Palace analysis is a choice of one
cell per axis over the same substrate---which is exactly why
\cref{ch:situating}'s grid has so few columns, and why this book could be written
at all.

\section{Two worked traces}
\label{sec:trace-worked}

We close the disassembly with two end-to-end traces small enough to follow by hand:
the driven filter at one frequency, and the electrostatic capacitor, each rendered
as a \emph{trace table} that names, at every stage, the component, the data shape,
and the chapter of origin.

\begin{workedexample}{The driven filter, one frequency, by data shape}{trace-driven}
We follow the data through a \emph{single} frequency $\omega = \omega_k$ of the
sweep---config to mesh to assemble to one GMRES solve to port projection to one
S-matrix entry---tracking the \emph{shape} of the value at each stage rather than
its numerical contents (the numerics are \cref{wex:opvary} and \cref{wex:sparam}
of \cref{ch:driven}; here we trace the \emph{pipeline}). Write $N$ for the number
of $\Hcurl$ degrees of freedom (millions, in practice) and $p=2$ for the number
of ports.

\medskip
\begin{center}\small
\setlength{\tabcolsep}{4pt}
\renewcommand{\arraystretch}{1.25}
\begin{tabular}{@{}ll>{\raggedright\arraybackslash}p{42mm}l@{}}
\toprule
Stage & Component & Data in $\to$ out & From \\
\midrule
parse    & \code{IoData} record      & file $\to$ \code{config}                  & \cref{ch:records} \\
mesh     & \code{build\_mesh}         & \code{config} $\to$ \code{Mesh}            & \cref{ch:functionspaces} \\
space    & \code{nd\_space}           & \code{Mesh} $\to$ $\Hcurl$ space (dim $N$) & \cref{ch:functionspaces} \\
assemble & \code{fe\_assemble}$^{\times3}$ & space $\to$ basis $\{\Kop,\Cop,\Mop\}$, each $N{\times}N$ & \cref{ch:assembly} \\
form     & \code{linear\_combination} & basis, $\omega_k$ $\to$ $\mat{A}(\omega_k)$ ($N{\times}N$, complex) & \cref{ch:combinators} \\
excite   & \code{excitation}          & \code{config}, $\omega_k$ $\to$ $\vb \in \Complex^N$ & \cref{ch:driven} \\
solve    & \code{ksp\_solve} (GMRES)  & $\mat{A}(\omega_k),\vb$ $\to$ $\vE_k \in \Complex^N$ & \cref{ch:gmres} \\
 \ \ precond. & V-cycle              & residual $\to$ correction ($\Complex^N$)   & \cref{ch:multigrid} \\
 \ \ smooth   & Chebyshev/Hiptmair    & per-level error $\to$ damped error         & \cref{ch:smoothers} \\
 \ \ apply    & matrix-free           & $\vw \in \Complex^N$ $\to$ $\mat{A}\vw$    & \cref{ch:matrixfree} \\
modes    & boundary-mode             & port faces $\to$ $\{\vs_1,\vs_2\}$         & \cref{ch:waveports} \\
reduce   & \code{sparameter\_reduce}  & $\vE_k, \{\vs_i\}$ $\to$ $S(\omega_k)$ ($p{\times}p$, complex) & \cref{ch:reductions} \\
\bottomrule
\end{tabular}
\end{center}

\medskip
Read the rightmost column top to bottom: it is a tour of the book. Read the
\emph{shape} column and watch the data collapse---from a file, to an $N$-dimensional
complex field, to a $2\times2$ complex matrix at this frequency. The big dimension
$N$ (the field) lives only \emph{inside} the solve; it is born at \code{nd\_space}
and dies at \code{sparameter\_reduce}, which projects it down to the handful of
port numbers the engineer wanted. One S-matrix entry---say $S_{21}(\omega_k)$, the
transmission---is a single complex number distilled from one $N$-dimensional GMRES
solve by one inner product against a port mode. Sweep this whole table across the
band and you have the plot of \cref{fig:trace-splot}; that sweep is the entire
driven driver.
\end{workedexample}

\begin{workedexample}{The electrostatic capacitor: the same skeleton, different cells}{trace-es}
Now run the identical lifecycle for the electrostatic capacitance extraction (the
clean contrast), and lay its trace table beside the driven one. The problem is $n$
conductors; write $N$ for the number of $\Hone$ degrees of freedom and $n$ for the
number of conductors.

\medskip
\begin{center}\small
\setlength{\tabcolsep}{4pt}
\renewcommand{\arraystretch}{1.25}
\begin{tabular}{@{}ll>{\raggedright\arraybackslash}p{30mm}>{\raggedright\arraybackslash}p{32mm}@{}}
\toprule
Stage & Component & Data in $\to$ out & Same as driven? \\
\midrule
parse    & \code{IoData} record      & file $\to$ \code{config}                  & \textbf{same} (\cref{ch:records}) \\
mesh     & \code{build\_mesh}         & \code{config} $\to$ \code{Mesh}            & \textbf{same} (\cref{ch:functionspaces}) \\
space    & \code{h1\_space}           & \code{Mesh} $\to$ $\Hone$ (dim $N$)        & \emph{swapped}: $\Hone$ not $\Hcurl$ \\
assemble & \code{fe\_assemble}         & space $\to$ $\Kop$ ($N{\times}N$, SPD)     & \textbf{same fold}, one term \\
excite   & \code{excitation}          & \code{config}, $j$ $\to$ $\vb_j \in \Real^N$ & per terminal (real) \\
solve    & \code{solve\_family} (CG)  & $\Kop, [\vb_j]$ $\to$ $[V_j] \in \Real^N$  & \emph{swapped}: corner~1, CG \\
 \ \ precond. & V-cycle              & residual $\to$ correction                   & \textbf{same} (\cref{ch:multigrid}) \\
 \ \ apply    & matrix-free           & $\vw$ $\to$ $\Kop\vw$                       & \textbf{same} (\cref{ch:matrixfree}) \\
reduce   & \code{gram\_reduce}        & $[V_j], \Kop$ $\to$ $\mat{C}$ ($n{\times}n$, sym.) & \emph{swapped}: Gram not projection \\
\bottomrule
\end{tabular}
\end{center}

\medskip
Set the two tables side by side and the book's whole thesis is visible in one
glance. The first two rows are \emph{identical}---same record, same mesh build. The
inner solver rows (V-cycle, matrix-free apply) are \emph{identical}---the same
numerical core. Exactly three rows differ: the space row ($\Hone$ for $\Hcurl$),
the solve row (\code{solve\_family}/CG/corner~1 for \code{frequency\_sweep}/GMRES/
corner~2), and the reduce row (\code{gram\_reduce} for \code{sparameter\_reduce}).
And one stage \emph{drops out} entirely: electrostatics needs no operator-forming
\code{linear\_combination} step, because its operator is fixed---which is precisely
the hoist of \cref{ch:electrostatics}. The electrostatic operator is also real and
SPD, so CG replaces GMRES and the V-cycle's smoother is a simpler real one; but the
\emph{stack is the same stack}. Two analyses, one machine, three swapped cells: the
trace tables prove it row by row.
\end{workedexample}

\begin{figure}[t]
\centering
\begin{tikzpicture}
\begin{axis}[
  width=0.84\linewidth, height=5.6cm,
  xlabel={frequency $f$ (GHz)}, ylabel={magnitude (dB)},
  xmin=3, xmax=7, ymin=-45, ymax=4,
  ytick={0,-10,-20,-30,-40}, xtick={3,4,5,6,7},
  legend pos=south east, legend cell align=left,
  grid=both, grid style={simGray!20},
  tick label style={font=\footnotesize}, label style={font=\footnotesize},
  legend style={font=\footnotesize},
]
% |S21| insertion loss: bandpass near 5 GHz
\addplot[very thick, simGreen, smooth, samples=120, domain=3:7]
  { -2 - 30*( (x-5)/1.0 )^2 / (1 + 0.55*((x-5))^2 ) - 0.5*(x-5)^2 };
\addlegendentry{$|S_{21}|$ (insertion loss)}
% |S11| return loss: deep dip at the match
\addplot[very thick, simRust, smooth, samples=160, domain=3:7]
  { -3 - 32*exp( -((x-5)^2)/0.06 ) + 2.0*((x-5))^2 };
\addlegendentry{$|S_{11}|$ (return loss)}
\draw[simBlue!50, densely dashed] (axis cs:4.4,4) -- (axis cs:4.4,-45);
\draw[simBlue!50, densely dashed] (axis cs:5.6,4) -- (axis cs:5.6,-45);
\node[font=\scriptsize, text=simBlue] at (axis cs:5.0,-43) {passband};
% mark one swept frequency = one full solve
\node[circle, fill=simInk, inner sep=1.1pt, label={[font=\scriptsize,text=simInk]above:{$\omega_k$}}]
  at (axis cs:4.7,-1.3) {};
\end{axis}
\end{tikzpicture}
\caption[The product: the swept S-parameter plot]{The output of the whole trace:
the scattering parameters of the two-port filter across the band. The insertion
loss $\abs{S_{21}}$ (green) sits near $0\,$dB in the passband around $5\,$GHz---signal
passes through---and rolls off outside it; the return loss $\abs{S_{11}}$ (rust)
dips sharply at the match. \emph{Every point on every curve is one complete pass of
the trace} of \cref{wex:trace-driven}: one $\mat{A}(\omega)$ formed, one GMRES solve
over the full multigrid stack, one port projection. The single marked point
$\omega_k$ is the frequency whose trace table we walked. Reading these two curves is
reading the device---and producing them is running the entire machine of
\cref{fig:trace-tree} once per frequency.}
\label{fig:trace-splot}
\end{figure}

\section{The bridge into Part~VI}
\label{sec:trace-bridge}

We have now done what Part~V set out to do. We met the six analyses, anchored them
all on electrostatics, and---in this chapter---ran one of them, the driven filter,
all the way through, naming every machine the book built and watching the whole
simulator collapse into a single composition (\cref{fig:trace-tree}). We saw, in two
trace tables, that a second analysis is the \emph{same} composition with three cells
swapped. The disassembly is complete: we know every part, and we know how they
compose into any driver.

The natural next question is the one Part~VI answers. If every analysis is a
composition of the same named pieces, then those pieces---taken together, with their
types and their bodies written out---are a small \emph{library}, and each driver is a
few lines of code over it. That is exactly the \emph{synthesized library} view. The
composition tree of \cref{fig:trace-tree} is not a metaphor; it is, almost literally,
the call graph of a real program, and Part~VI writes that program down. \Cref{ch:synthlib}
lays out the library's architecture---the handful of modules (types, iteration,
data-algebra, coordination, drivers) that hold the machines we have been naming---and
\cref{ch:composing} composes a driver end to end in code, the rendered form of the
very tree we just expanded.

\begin{keyidea}
\textbf{You have now seen the whole machine.} One simulation is one composition of
the same handful of machines---a mesh build, an assembly fold, a solve corner with
its nested Krylov/multigrid/matrix-free stack, and a reduction---and a different
analysis swaps only a few of those boxes. Part~V disassembled the simulator into
those boxes and showed they compose; Part~VI reassembles them into a single small
library and renders the whole simulator as a few pages of code. The trace tree of
this chapter \emph{is} that program, drawn rather than typed.
\end{keyidea}

\summarysection
This capstone traces one complete simulation---the driven frequency sweep of a
two-port filter---from the configuration file to the S-parameter plot, naming at
every stage the exact machine, and chapter, that does the work. Walking the
lifecycle of \cref{ch:lifecycle}: \emph{parse} produces the read-only \code{IoData}
record (\cref{ch:records}, \cref{ch:strata}); \emph{build mesh} produces the
\code{Mesh} (\cref{ch:functionspaces}); \emph{dispatch} selects the driven driver at
the specialization seam (\cref{ch:lifecycle}); \emph{assemble} folds three weak-form
term lists (\cref{ch:weak}) into the fixed basis $\{\Kop,\Cop,\Mop\}$ by
\code{fe\_assemble} (\cref{ch:assembly}), with matrix-free applies
(\cref{ch:matrixfree}) and eliminated boundary conditions (\cref{ch:bc});
\emph{solve} is the operator-varying map \code{frequency\_sweep} (\cref{ch:driven}),
each frequency a preconditioned GMRES (\cref{ch:gmres}) over a multigrid V-cycle
(\cref{ch:multigrid}) whose smoother is Chebyshev/Hiptmair (\cref{ch:smoothers}) and
whose every apply is matrix-free, the state threaded by the solve monad
(\cref{ch:monad}); \emph{reduce} is the port-projection \code{sparameter\_reduce}
(\cref{ch:reductions}) over port modes borrowed from the boundary-mode analysis
(\cref{ch:waveports}). The whole driver is one composition, \texttt{driven $=$
sparameter\_reduce $\circ$ frequency\_sweep $\circ$ fe\_assemble}, which expands into
a single finite tree whose every leaf is an earlier chapter (\cref{fig:trace-tree}).
Set beside it, the electrostatic capacitance trace is the \emph{same} machine---same
record, mesh, assembly fold, inner solver stack---with exactly three cells swapped:
the function space ($\Hone$/$\Hcurl$), the solve corner (fixed-operator map/CG versus
operator-varying map/GMRES), and the reduction (symmetric Gram versus
port-projection). One run touches one solve corner and one reduction shape on the
grid of \cref{ch:situating}: the simulator is a few corners $\times$ a few reductions
over one shared field theory and calculus. Part~VI (\cref{ch:synthlib},
\cref{ch:composing}) renders that composition as code.

\furtherreading
The scattering-parameter physics this trace computes---ports, impedance
normalization, return and insertion loss, the network-analyzer measurement---is the
foundation of microwave engineering, and Pozar's text~\citep{pozar2011} is the
standard treatment; the finite-element formulation of the time-harmonic curl--curl
system on $\Hcurl$, the wave-port modes, and the modal projection are developed at
length by Jin~\citep{jin2014}. The numerical engine the trace nests four
deep---GMRES on the complex indefinite operator, the multigrid preconditioner with
its smoother, and the reuse economics across the family---is the subject of Saad's
account of iterative methods~\citep{saad2003}, on which \cref{ch:gmres},
\cref{ch:multigrid}, and \cref{ch:smoothers} all draw. The individual machines named
in this trace are each developed in their own chapters, cited inline; this chapter
adds no new content, only the composition---which Part~VI (\cref{ch:synthlib},
\cref{ch:composing}) then renders as a working library.

\exercisessection
\begin{enumerate}[nosep]
  \item \textbf{Name the machine at each stage.} Without looking back at
  \cref{fig:trace-pipeline}, list the lifecycle stages of the driven trace in order
  (parse, mesh, dispatch, assemble, solve, reduce, output) and name the framework
  component and the chapter responsible for each. Which single stage consumes a
  product computed by a \emph{different} analysis, and which analysis?
  \item \textbf{Expand one branch.} Take the \code{frequency\_sweep} box of the
  composition tree (\cref{fig:trace-tree}) and write out, as a nested list, every
  machine that engages inside a single \code{ksp\_solve}---from the outer Krylov
  method down to the innermost operation. At what operation does \emph{every} branch
  of the tree bottom out, and why is that the same operation for GMRES, the V-cycle,
  and the smoother alike?
  \item \textbf{The one line that swaps the analysis.} The driven and electrostatic
  traces share most of their stages. Identify the single operation whose
  \emph{position} (inside the solve loop versus hoisted outside it) flips the run
  from corner~1 to corner~2, and explain why electrostatics can hoist it while the
  driven sweep cannot. (Cross-check against \cref{ch:driven}.)
  \item \textbf{Trace the shape.} Following the trace table of
  \cref{wex:trace-driven}, state the data shape carried out of each stage for the
  driven filter, from the config file to one S-matrix entry. At which stage is the
  large dimension $N$ (the field) born, and at which does it die? What is the shape
  of the final product at one frequency, and how many ports does it have on a side?
  \item \textbf{Two analyses, three swaps.} Using the two trace tables
  (\cref{wex:trace-driven}, \cref{wex:trace-es}), list exactly which rows are
  identical between the driven and electrostatic traces and which three are swapped.
  Which stage present in the driven trace \emph{drops out} entirely for
  electrostatics, and why? (Hint: what does a fixed operator not need to do each
  iteration?)
  \item \textbf{Count the footprint.} Of the four solve corners and three reduction
  shapes of \cref{ch:situating}, how many of each did this one driven run use? Name
  them. The run also \emph{borrowed} a product from a corner it did not itself
  occupy---which corner, and what did it borrow? Repeat the count for the
  electrostatic trace.
  \item \textbf{A third analysis on the same machine.} Pick the \emph{transient}
  analysis (\cref{ch:transient}). Predict, by analogy with the two traces of this
  chapter, which cells of its trace table would be \emph{shared} with the driven
  filter (record, mesh, assembly fold, inner stack) and which would be
  \emph{swapped}. Which solve corner does it occupy, and why is its solve a
  \emph{fold} rather than a \emph{map}?
  \item \textbf{From tree to code.} The composition tree of \cref{fig:trace-tree} is
  described in the bridge (\S\ref{sec:trace-bridge}) as ``almost literally the call
  graph of a real program.'' In two or three sentences, explain what would have to
  be written down to turn the tree into running code: which leaves are already-built
  library functions, which boxes are the composition itself, and which single piece
  (named in \cref{ch:waveports}) must be supplied as an input. (Part~VI,
  \cref{ch:synthlib} and \cref{ch:composing}, does exactly this.)
\end{enumerate}
